% --------------------------------- Preamble -----------------------------------
\documentclass[11pt, letterpaper, twocolumn]{article}

\usepackage[T1]{fontenc} % use 8-bit fonts
\usepackage[margin=0.75in, columnsep=0.2in, % 3.4 in columns
    headsep=0.28in, footskip=0.42in]{geometry} % paper sizes
\usepackage{amsmath, amsfonts, amssymb} % improved math
\usepackage{bm} % bold math (must be after amsmath)
\usepackage{upquote} % upright single quotes in verbatim
\usepackage{pgfplots} % includes tikz, xcolor, graphicx
\usepackage[colorlinks=true, allcolors=gray]{hyperref}

\raggedbottom % do not force stretching text to fill the page
\setlength{\emergencystretch}{\hsize} % avoid overfull lines
\overfullrule=1mm % black band to show overfull lines
\pdfsuppresswarningpagegroup=1 % remove pdf image warnings
\frenchspacing % consistent spacing between words
\newcommand{\ul}[1] % vectors
    {{}\mkern1mu\underline{\mkern-1mu#1\mkern-1mu}\mkern1mu}
\DeclareSymbolFont{euler}{U}{eur}{m}{n} % symbol font
\DeclareMathSymbol{\PI}{\mathord}{euler}{25} % up pi
\pgfplotsset{compat=newest} % use newest pgfplot version
\newcommand{\EE}[2]{\ensuremath{{{#1}\!\times\!10^{#2}}}}
% ------------------------------------------------------------------------------

\title{Discrete Bayes Filter for Circular Robot Localization}
\author{}
\date{}
\begin{document}
\maketitle

\section*{Description}

A toy boat randomly moves about the edge of a circularly looped pool with 50
colored tiles on the floor. There are 5 different colors. \\
\includegraphics[width=\linewidth]{../figures/map.png}
The only sensor the toy has detects the color of the tile immediately beneath
the boat. Because of the water, the sensor does not always detect the correct
color. The toy boat also has a map of the colored tiles with which to compare
the sensor readings. The wind tends to move the boat around the pool in one
direction more than another. With knowledge of the map, the sensor readings, the
PMF of the sensor, and the PMF of the dynamics, the boat can figure out where it
is in the pool.

\section*{Bayes Filter}
The robot's position is estimated using a discrete Bayes filter. The belief over
the state, representated as a probability mass function $p(x)$, is updated
recursively through two steps: \\

\textbf{Measurement Update:} The prior belief is updated using the likelihood
of the observed measurement:
\[
p(x) \propto p(z{\mid}x) p(x)
\]

\textbf{Time Propagation:} The updated belief is propagated forward using the 
Markov transition matrix:
\[
p(x_{k+1}) = Mp(x_k)
\]

\section*{Performance Metrics}

Three metrics are used to evaluate the filter:

\textbf{Accuracy:} The fraction of time steps where the estimated position,
given by the maximum of $p(x)$, matches the true position.

\textbf{Validity:} A measure of how well the probability distribution represents
the true state. When the estimate matches the truth, the contribution is weighted
by the inverse of the estimate probability. A valid estimator produces a value
close to 1.

\textbf{Goodness:} The fraction of time steps where the maximum probability exceeds
a threshold of 0.5, indicating high confidence in the estimate.

\section*{Monte Carlo Simulation}

To evaluate the performance of the estimator, all experiments are conducted using
Monte Carlo simulations with 100 independent runs. For each run, a new trajectory
and set of measurements are generated, and the resulting accuracy, validity, and
goodness metrics averaged across all runs.

\subsection*{Nominal Case}

The filter is first evaluated using matched parameters for both data generation and
estimation, with sensor reliability $p_c = 0.8$ and motion model $p_m = [0.2, 0.3, 0.5]$.
This case serves as a baseline and is expected to produce a validity value close to 1,
indicating that the estimated probability mass function accurately represents the 
true system behavior.

\subsection*{Sensor Model Mismatch}

To study the effect of model mismatch, the true sensor reliability is fixed at $p_c = 0.8$,
while the estimator assumes different values.

First, the estimator assumes $p_c = 0.6$, underestimating the sensor accuracy. Next,
the estimator assumes $p_c = 0.99$, overestimating the sensor accuracy.

\subsection*{Parameter Sweep}

A broader analysis is conducted by varying both the sensor reliability and 
motion model. Four motion models are considered:

\begin{center}
    \begin{tabular}{c c}
        \hline
        Set & $p_m$ \\
        \hline
        1 & [0.3, 0.3, 0.4] \\
        2 & [0.2, 0.2, 0.6] \\
        3 & [0.1, 0.1, 0.8] \\
        4 & [0.1, 0.8, 0.1] \\
        \hline
    \end{tabular}
\end{center}

For each motion model, the sensor reliability $p_c$ is varied from 0.6 to
0.95 in increments of 0.05. This results in 32 total scenarios. For each
scenario, the accuracy metric is computed as the average over 100 Monte
Carlo runs.

The resulting accuracy curves illustrate how estimation performance depends
on both sensor quality and system dynamics.

\section*{Results}

\subsection*{Nominal Case}

Figure~\ref{fig:nominal} shows the results when both generation and
estimation use the same \texttt{pc} value.

\begin{figure}[h]
    \centering
    \includegraphics[width=\linewidth]{../figures/fig_results_80.pdf}
    \caption{Validity: 1.001     Goodness: 0.59     Accuracy: 0.5522}
    \label{fig:nominal}
\end{figure}

\pagebreak

\subsection*{Sensor Model Mismatch}

Figure~\ref{fig:underestimator} shows the results when only estimation
uses a \texttt{pc} value of 0.6.

\begin{figure}[h]
    \centering
    \includegraphics[width=\linewidth]{../figures/fig_results_60.pdf}
    \caption{Validity: 1.3344     Goodness: 0.31     Accuracy: 0.5413}
    \label{fig:underestimator}
\end{figure}

The validity and goodness scores differed significantly in this scenario.
Since the \texttt{pc} value was less for our estimator than reality, the estimator
was underconfident in our sensor. This lead to a validity higher than 1 because our 
estimator believed our sensor to be less accurate than it truly was. This also
decreased our goodness score because our peak probabilities were lower due to the 
decreased belief in the accuracy of our sensor. \\

Figure~\ref{fig:overestimator} shows the results when only estimation
uses a \texttt{pc} value of 0.99.

\begin{figure}[h]
    \centering
    \includegraphics[width=\linewidth]{../figures/fig_results_99.pdf}
    \caption{Validity: 0.7776     Goodness: 0.76     Accuracy: 0.5413}
    \label{fig:overestimator}
\end{figure}

The validity and goodness scores again differed in this scenario. Since the 
\texttt{pc} value was greater for our estimator than reality, the estimator was
overconfident in our sensor. This lead to a validity lower than 1 because our estimator
believed our sensor to be more accuracte than it truly was. This also increased our
goodness score because our peak probabilities were higher due to the increased belief
in the accuracy of our sensor.

\subsection*{Parameter Sweep}

Figure~\ref{fig:accuracy} shows, in a single figure, the \texttt{accuracy}
scores for each of the \texttt{pm} sets, labeling each curve ``Set \{number\}''.

\begin{figure}[h]
    \centering
    \includegraphics[width=\linewidth]{../figures/fig_accuracy.pdf}
    \caption{Accuracies of each motion model over 100 Monte Carlo runs}
    \label{fig:accuracy}
\end{figure}

The first observation worth noting is the most obvious and self-explanatory.
As the sensor accuracy increases, the estimator accuracy increases. This makes
sense since a good estimator would be more confident in higher accuracy measurements.

The second observation of importance is the general rule of increased estimator accuracy
for larger dynamics probabilities. What I mean by this is that the most accurate
estimator was the one for the dynamics PMF: [0.1, 0.1, 0.8]. Compared to the estimator
for the dynamics PMF: [0.3, 0.3, 0.4], the accuracy was much greater for the one
with the greatest single probability in the PMF. This makes sense because if the 
dynamics say the robot will move left with probability 0.8, the estimator will
become very confident that the robot will continue to move forward. However, if
the robot moves closer to a uniform distribution, it is much harder for the estimator
to predict where the robot will go next, leading to a lower estimator accuracy.

The last major observation I noticed was quite unintuitive. The outlier for the 
rule I just mentioned above was the PMF: [0.1, 0.8, 0.1]. Even though it had the highest
single probability, the estimator accuracy was actually the lowest compared to the
other PMFs. This is mainly due to the fact that the dynamics of this robot mean
it will most likely stay put, compared to the robot of [0.1, 0.1, 0.8] being most
likely to move right (and explore). This hurts the estimators accuracy for two reasons:
The first being that each step provides very little new information when the robot is
mostly dormant. If the robot is staying put over many steps, the filter will only 
become slightly more confident in its estimation. The second reason is that the filter
will have significantly more trouble recovering from errors. If the dynamics say the robot
is much more likely to stay stationary and the measurement has lead our estimator to 
believe we are in the wrong position, our estimator may be stuck in that position
for a longer period of time. Whereas the estimator for the dynamics of the exploring
robot will get new measurements at a much greater frequency, giving the filter 
more information to process and make a good estimate.

\section*{Discussion}

The results demonstrate that accurate localization depends on both the quality of 
the sensor model and dynamics of the system. When the assumed sensor model matches
the true system, the filter produces valid probability estimates. However, mismatches
in the assumed sensor reliability lead to degrade performance, particularly in 
the validity metric.

Higher sensor accuracy generally improves estimation performance, while motion
models with stronger directional bias lead to more predictable and accurate
localization. These results highlight the importance of correctly modeling 
both system dynamics and measurement uncertainty in Bayesian estimation.

\end{document}